\documentclass[10pt,a4paper]{article}

\usepackage[T1]{fontenc}
\usepackage[utf8]{inputenc}
\usepackage{lmodern}
\usepackage[margin=0.7in]{geometry}
\usepackage{xcolor}
\usepackage[urlbordercolor=cyan, linkbordercolor=cyan]{hyperref}
\usepackage{titlesec}
\usepackage{enumitem}
\usepackage{fontawesome5}

\titleformat{\section}
  {\normalsize\scshape\bfseries}{}{0em}{}[\titlerule\vspace{2pt}]
\titlespacing*{\section}{0pt}{8pt}{2pt}

\setlength{\parindent}{0pt}
\pagenumbering{gobble}

% Custom bullets
\setlist[itemize,1]{label=$\bullet$,leftmargin=1.2em,topsep=2pt,itemsep=2pt,parsep=0pt}
\setlist[itemize,2]{label=$\circ$,leftmargin=1.2em,topsep=1pt,itemsep=1pt,parsep=0pt}

% Entry command: name + right-aligned date, role + right-aligned location
\newcommand{\entry}[4]{%
  \textbf{#1} \hfill \textit{#2} \\
  \textit{#3} \hfill #4
}

\begin{document}

\begin{center}
{\LARGE \textbf{Seonghyeon Yun}} \\[2pt]
% \href{mailto:tjdgus4384@snu.ac.kr}{tjdgus4384@snu.ac.kr} \quad $\vert$ \quad
\href{https://github.com/tjdgus4384}{\faGithub\ tjdgus4384} \quad $\vert$ \quad
\href{https://www.linkedin.com/in/seonghyeon-yun-0096142a0}{\faLinkedin\ Seonghyeon Yun} \\
Seoul, South Korea
\end{center}

\section*{Education}
\begin{itemize}
  \item \entry{Seoul National University}{Mar.~2024 -- Feb.~2028 (expected)}{B.S.~in Computer Science and Engineering}{Seoul, Korea}
  \item \entry{Seoul Science High School}{Mar.~2021 -- Feb.~2024}{High School diploma}{Seoul, Korea}
\end{itemize}

\section*{Research Experience}
\begin{itemize}
  \item \entry{\href{https://jaesik.info/lab}{Visual \& Geometric Intelligence Lab}, Seoul National University}{Mar.~2026 -- Present}{Undergraduate Researcher (Advisor: \href{https://jaesik.info}{Prof.~Jaesik Park})}{Seoul, Korea}
  \begin{itemize}
    \item Research on 3D reconstruction, robot manipulation, and grasping.
  \end{itemize}

  \item \entry{\href{https://www.telepix.net}{TelePIX Co., Ltd.}}{Jan.~-- Feb.~2026}{AI Research Team, Cosmo Pioneers Fellowship}{Seoul, Korea}
  \begin{itemize}
    \item Selected as 1st-cohort fellow ($\sim$1 in 75 acceptance rate).
    \item 3D reconstruction from sparse satellite imagery using Gaussian Splatting.
  \end{itemize}
\end{itemize}

\section*{Publications}
\begin{itemize}
  \item \textbf{NL-ArgCheck: Natural-Language Argument Verification}. Self-driven research, Ongoing.
  \item \textbf{GeoGS: Geospatial Gaussian Splatting for Robust 3D Reconstruction from Sparse Satellite \mbox{Imagery}}. Han-Gyeol Kim$^*$, \textbf{Seonghyeon Yun}$^*$, JaeWan Park, Darongsae Kwon. \emph{CVPR 2026 EarthVision Workshop} \textbf{(Best Paper Award)}. ($^*$ equal contribution) \href{https://github.com/tjdgus4384/GeoGS}{[code]}
\end{itemize}

\section*{Honors and Awards}
\begin{itemize}
  \item \textbf{National Science and Engineering Scholarship} \hfill \textit{2025 -- Present} \\
        \emph{Ministry of Science and ICT, Republic of Korea}
        \begin{itemize}
          \item Elected as a 12th-cohort cadet of Research Officers for National Defense. $\sim$25 selected nationwide per year.
          \item $\sim$USD 4,000 per semester for tuition and stipends.
        \end{itemize}
  \item \textbf{Excellence Award, ``2050 \emph{Mirae-in}: AI and Our Lives'' Scenario Contest} \hfill \textit{2024} \\
        \emph{Art Center Nabi \& Seoul National University AI Institute, AI ELSI Center} (team award)
  \item \textbf{Outstanding Research Award (1st in Information Science Track)} \hfill \textit{2022} \\
        \emph{Creative Interdisciplinary Research, Seoul Science High School} -- Project: \emph{Multi-agent Reinforcement Learning for Grid Soccer using DQN}
  \item \textbf{Certificate of Completion, International Olympiad in Informatics Summer School} \hfill \textit{2021} \\
        \emph{Korea Information Science Society}
  \item \textbf{Silver Award, National Finals, Korea Olympiad in Informatics (KOI)} \hfill \textit{2019, 2020}
\end{itemize}

\section*{Projects}
\begin{itemize}
  \item \textbf{Aave V3 \& Morpho Blue Liquidation MEV Bot (Base)} \\
        Same-block, oracle-driven liquidation bot for the Aave V3 and Morpho Blue lending protocols on Base L2, executed atomically via Balancer flash loans and a Uniswap V3 swap.
        \begin{itemize}
          \item \emph{Off-chain Chainlink \texttt{transmit()} predictor}: aggregated CEX mid (OKX + Coinbase WebSockets) drives a deviation/heartbeat trigger that anticipates on-chain oracle updates with $\sim$46\,s median lead, validated on a 30-day historical replay --- the only viable predict signal on Base, which exposes no public mempool.
          \item \emph{On-chain forensics \& strategy correction}: established empirically that Base block ordering is concatenated flashblocks ($\sim$250\,ms sub-blocks), not a single global priority-fee sort, by analyzing per-tx ordering vs.\ effective priority on historical same-block races; refit the bidding model and priority ceiling accordingly.
          \item \emph{Realistic profit backtest}: 121 same-block races / 7 days, \$76k gross weekly spread; competition-aware net (winner-identity model + crash-slippage stress + Uniswap V3 \texttt{QuoterV2} swap impact) projects \$53--74k/wk under static-incumbent assumptions.
          \item \emph{Production hardening}: Foundry-tested Solidity \texttt{LiquidatorV1} (on-chain \texttt{minProfitWei} guard, owner-pausable kill switch); Python/\,asyncio engine with supervised workers, per-opportunity revert circuit breaker, multi-RPC failover, and a \texttt{setsid}-detached supervisor (auto-restart + cron \texttt{@reboot}) for unattended live operation.
        \end{itemize}
  \item \textbf{Parallel ARC-AGI Solver via Transductive LLMs and Rule-Based CPU Search} \\
        Team project, SNU Basics of Deep Learning (Spring 2025). Solving an easier subset of the \href{https://arcprize.org}{ARC-AGI} benchmark. Transductive Llama-3.1-8B (LoRA + Test-Time Training) running concurrently with rule-based hard-coded solutions.
  \item \textbf{ClassBridge} \hfill \href{https://github.com/tjdgus4384/ClassBridge}{[\faGithub]} \\
        Real-time anonymous classroom communication tool. Instructors monitor student comprehension status via a desktop widget; students join via QR scan to signal understanding or submit anonymous questions.
        \begin{itemize}
          \item \textit{Stack:} TypeScript, Next.js, Tailwind CSS, Node.js, Socket.IO, Electron, Fly.io.
        \end{itemize}
\end{itemize}

\section*{Skills}
\begin{itemize}
  \item \textbf{Programming Languages:} Python, C++, TypeScript
  \item \textbf{Frameworks/Tools:} PyTorch, Git, \LaTeX, Next.js, Socket.IO, Electron
  \item \textbf{Areas:} 3D Computer Vision, Gaussian Splatting / NeRF, Robot Learning, Machine Learning Theory
  \item \textbf{Languages:} Korean (native), English (fluent)
\end{itemize}

\end{document}
